\documentclass[12pt, a4paper]{article}

% --- Pakete ---
\usepackage[utf8]{inputenc}
\usepackage[T1]{fontenc}
\usepackage[ngerman]{babel}
\usepackage{csquotes}
\usepackage{microtype}
\usepackage{geometry}
\geometry{margin=2.5cm}

\usepackage{amsmath}
\usepackage{booktabs}
\usepackage{graphicx}
\usepackage{xcolor}
\usepackage{hyperref}
\hypersetup{
    colorlinks = true,
    linkcolor  = black,
    citecolor  = black,
    urlcolor   = blue,
}

% Literatur
\usepackage[backend=biber, style=authoryear, sorting=nyt]{biblatex}
\addbibresource{literatur.bib}

% --- Metadaten ---
\title{%
  \textbf{Nachhaltiges Prescriptive Process Monitoring}\\[0.4em]
  \large Ein graphbasierter Ansatz zur ökologischen Prozesssteuerung
  mittels erweiterter Eventlogs und dem PIXIE-Algorithmus
}
\author{Johanna Prinz}
\date{\today}

% ============================================================
\begin{document}
\maketitle
\tableofcontents
\newpage

% ------------------------------------------------------------
\section{Einleitung}
\label{sec:einleitung}

Klassische Prozess-Eventlogs erfassen Aktivitäten, Zeitstempel und Kosten~--
ökologische Auswirkungen bleiben dabei unsichtbar.
Klimaziele und regulatorische Anforderungen erfordern jedoch,
dass Unternehmen Nachhaltigkeitskennzahlen nicht erst \emph{nach}
Prozessabschluss in Post-Mortem-Analysen auswerten,
sondern \emph{während} der Prozessausführung aktiv steuern können.

Die vorliegende Arbeit schließt diese Forschungslücke,
indem sie erweiterte Eventlogs, Predictive Monitoring
und einen graphbasierten Empfehlungsalgorithmus zu einem
einheitlichen \emph{nachhaltigen Prescriptive-Monitoring-Ansatz} verbindet.

% ------------------------------------------------------------
\section{Erweiterung klassischer Eventlogs um Nachhaltigkeitsattribute}
\label{sec:eventlog}

Ein klassisches Eventlog enthält je Ereignis mindestens:
\begin{itemize}
  \item \texttt{case\_id} -- Prozessinstanz,
  \item \texttt{activity} -- ausgeführte Aktivität,
  \item \texttt{timestamp} -- Ausführungszeitpunkt,
  \item \texttt{cost} -- anfallende Kosten.
\end{itemize}

Im vorgeschlagenen Ansatz wird jedes Ereignis um
\emph{quantifizierbare Nachhaltigkeitswerte} ergänzt:

\begin{equation}
  e = \bigl(\text{case\_id},\;\text{activity},\;\text{timestamp},\;
            \underbrace{\text{cost},\;\text{co}_2,\;\text{energy}}_{\text{neu}}\bigr)
\end{equation}

\noindent
Dabei bezeichnen:
\begin{itemize}
  \item $\text{co}_2$ -- \COO-Emissionen der Aktivität in~kg,
  \item $\text{energy}$ -- Energieverbrauch in~kWh,
  \item $\text{cost}$ -- monetäre Kosten in~\euro{}.
\end{itemize}

Ökologische Faktoren sind damit direkt im Eventlog
als messbare Attribute verankert und stehen für eine
Echtzeit-Analyse zur Verfügung.

% ------------------------------------------------------------
\section{Graphmodellierung des laufenden Prozesses}
\label{sec:graph}

Der laufende Prozess wird als \textbf{bipartiter gerichteter Graph}
$G = (V, E)$ modelliert:
\begin{itemize}
  \item Linke Knotenmenge~$S$: aktuelle Prozesszustände (beobachtete Aktivitäten),
  \item Rechte Knotenmenge~$A$: mögliche nächste Aktivitäten,
  \item Kante $(s, a) \in E$: historisch beobachteter Übergang mit den Attributen
        Häufigkeit~$w$, $\overline{\text{co}_2}$, $\overline{\text{cost}}$,
        $\overline{\text{energy}}$ (jeweilige Mittelwerte aus dem Eventlog).
\end{itemize}

Aus dem historischen Eventlog wird der Graph automatisch aufgebaut:
Für jede Prozessinstanz werden aufeinanderfolgende Aktivitäten
als gerichtete Kante eingetragen und deren Attributwerte aggregiert.

% ------------------------------------------------------------
\section{PIXIE-Algorithmus: Biased Random Walk zur Pfadbewertung}
\label{sec:pixie}

\subsection{Grundprinzip}

Der PIXIE-Algorithmus~\parencite{pixie2018} traversiert den Prozessgraphen
mittels eines \emph{biased Random Walk}:
Ausgehend vom aktuellen Prozesszustand~$s$ wird in jedem Schritt
ein Nachfolger~$a$ probabilistisch gewählt,
wobei die Übergangswahrscheinlichkeit durch eine
\textbf{multidimensionale Bewertungsfunktion} gesteuert wird.

\subsection{Bewertungsfunktion}

Für jeden möglichen Übergang $(s, a)$ wird ein Score berechnet:

\begin{equation}
  \label{eq:score}
  \text{score}(s,a)
    = \alpha \cdot \hat{w}(s,a)
    + \beta  \cdot \bigl(1 - \hat{\text{cost}}(s,a)\bigr)
    + \gamma \cdot \bigl(1 - \widehat{\text{co}_2}(s,a)\bigr)
    + \delta \cdot \bigl(1 - \hat{\text{energy}}(s,a)\bigr)
\end{equation}

\noindent
wobei $\hat{\cdot}$ eine Min-Max-Normalisierung auf $[0,1]$ bezeichnet
und die Parameter $\alpha, \beta, \gamma, \delta \geq 0$
die Unternehmensprioritäten widerspiegeln:

\begin{table}[h]
  \centering
  \caption{Bedeutung der Gewichtungsparameter}
  \label{tab:parameter}
  \begin{tabular}{clll}
    \toprule
    Parameter & Dimension & Hoher Wert bedeutet & Beispielszenario \\
    \midrule
    $\alpha$ & Historische Häufigkeit  & Bewährte Pfade bevorzugen & Produktionsleiter \\
    $\beta$  & Niedrige Kosten         & Kostenoptimierung          & Controller \\
    $\gamma$ & Niedriger \COO-Ausstoß  & Klimaschutz                & Nachhaltigkeitsbeauftragter \\
    $\delta$ & Niedriger Energiebedarf & Energieeffizienz           & Energiemanager \\
    \bottomrule
  \end{tabular}
\end{table}

\subsection{Übergangswahrscheinlichkeiten}

Die Wahrscheinlichkeit, von Zustand~$s$ zur Aktivität~$a$ überzugehen, beträgt:

\begin{equation}
  p(a \mid s) = \frac{\text{score}(s,a)}{\sum_{a' \in \mathcal{N}(s)} \text{score}(s,a')}
\end{equation}

Mit Restart-Wahrscheinlichkeit~$r$ kehrt der Walk
bei jedem Schritt mit Wahrscheinlichkeit~$r$ zum Startknoten zurück,
was den Einfluss weit entfernter Knoten dämpft.

% ------------------------------------------------------------
\section{Prescriptive Monitoring: Proaktive Handlungsempfehlung}
\label{sec:prescriptive}

Nach $n$ Schritten des Random Walk wird die Besuchshäufigkeit
jedes Zielknotens ausgewertet.
Die Aktivität mit der höchsten Besuchsrate erhält die
\textbf{Empfehlung als nächster Prozessschritt}:

\begin{equation}
  a^* = \arg\max_{a} \;\text{visits}(a)
\end{equation}

Dieses Prescriptive Monitoring wirkt \emph{proaktiv}:
Es erkennt frühzeitig, ob ein laufender Prozess
auf einen ökologisch ungünstigen Pfad zusteuert --
noch bevor der Schaden eingetreten ist.

\subsection{Trade-off-Transparenz}

Der Algorithmus macht Zielkonflikte explizit sichtbar.
Durch systematische Variation von~$\gamma$ lässt sich darstellen:

\begin{quote}
  \emph{,,Pfad~B dauert 15\,\% länger, spart aber 40\,\% \COO~gegenüber Pfad~A.''}
\end{quote}

Entscheidungsträger erhalten damit eine transparente Grundlage,
um Schnelligkeit, Wirtschaftlichkeit und Ökologie
entsprechend der aktuellen Unternehmensstrategie abzuwägen.

% ------------------------------------------------------------
\section{Prototypische Implementierung}
\label{sec:implementierung}

Der Ansatz wurde prototypisch in Python implementiert.
Tabelle~\ref{tab:szenarien} zeigt die Empfehlungen
für drei Personas bei identischem Startknoten (\emph{Qualitätskontrolle}).

\begin{table}[h]
  \centering
  \caption{Szenariobasierte Empfehlungen des PIXIE-Algorithmus}
  \label{tab:szenarien}
  \begin{tabular}{lcccc}
    \toprule
    Persona & $\alpha$ & $\beta$ & $\gamma$ & Empfehlung \\
    \midrule
    Produktionsleiter (schnell)   & 0{,}6 & 0{,}2 & 0{,}1 & Lackieren Ofen \\
    Controller (kostenoptimiert)  & 0{,}2 & 0{,}6 & 0{,}1 & Lackieren UV-Aushärtung \\
    Nachhaltigkeitsbeauftragter   & 0{,}1 & 0{,}1 & 0{,}6 & Lackieren UV-Aushärtung \\
    \bottomrule
  \end{tabular}
\end{table}

% ------------------------------------------------------------
\section{Forschungsbeitrag und Abgrenzung}
\label{sec:beitrag}

Die Arbeit schließt eine bestehende Forschungslücke:
Erstmals werden drei bisher getrennte Konzepte zu einem
kohärenten Ansatz kombiniert:

\begin{enumerate}
  \item \textbf{Erweiterte Eventlogs} mit quantifizierbaren Nachhaltigkeitsattributen,
  \item \textbf{Predictive Monitoring} zur Früherkennung ökologisch ungünstiger Pfade,
  \item \textbf{Graphbasierter Empfehlungsalgorithmus} (PIXIE)
        für eine flexible, priorisierungsgesteuerte Handlungsempfehlung.
\end{enumerate}

Im Gegensatz zu klassischen Post-Mortem-Analysen
ermöglicht der Ansatz eine \emph{Echtzeit-Intervention}
während der Prozessausführung.

% ------------------------------------------------------------
\section{Fazit}
\label{sec:fazit}

Der vorgestellte Ansatz verbindet Process Mining,
nachhaltige Unternehmensführung und algorithmisches Empfehlungsmanagement
zu einem praxistauglichen Werkzeug.
Die flexible Gewichtungsfunktion erlaubt es,
denselben Algorithmus für unterschiedliche Unternehmensstrategien einzusetzen --
von maximaler Kosteneffizienz bis zur klimaneutralen Produktion.

% ------------------------------------------------------------
\printbibliography[heading=bibintoc]

\end{document}
